\chapter[Laminations, entropy, and holomorphic motions]{Source digests V: laminations, core entropy, and holomorphic motions}
\label{ch:digests-laminations-entropy}

This chapter digests three sources that the combinatorial and phase--parameter
approaches use: Dudko's \emph{Matings with laminations}~\cite{DudkoMatings}, the
Dudko--Schleicher paper on core entropy~\cite{DudkoSchleicher}, and Słodkowski's
theorem on holomorphic motions~\cite{Slodkowski}.  Statements are those of the sources,
up to notation; numbering is that of the source.

\section{Matings with laminations}
\label{sec:digest-matings}

\subsection*{Definitions}

\begin{definition}[Mandelbrot set and the rational slice $V_2$]
\label{def:matings-mandelbrot}
The Mandelbrot set is $\M=\{c:\ J(z^2+c)\text{ is connected}\}$.  The space $V_2$ is the
set of quadratic rational maps with a superattracting cycle of period two; each element
is uniquely of the form $g_a(z)=a/(z^2+2)$ with $a\in\C\setminus\{0\}$.  The nonescaping
locus $M_2\subset V_2$ consists of the parameters for which the second critical point is not
in the attracting basin of the superattracting $2$-cycle.
\end{definition}

\begin{definition}[Basilica lamination]
\label{def:basilica-lamination}
A collection $L_B$ of leaves in $\overline{\mathbb D}$ satisfies:
(1) each leaf joins points $\frac{a}{3\cdot2^n},\frac{b}{3\cdot2^n}\in S^1$, with $a,b$ coprime to $6$,
and each such point lies on exactly one leaf; (2) $\{\tfrac13,\tfrac23\}$ and $\{\tfrac56,\tfrac16\}$
are joined; (3) if $\{a,b\}\neq\{\tfrac13,\tfrac23\}$ is joined, the pairs among $\{a,b,\tfrac a2,\tfrac b2\}$
lying in the same component of $S^1\setminus\{\tfrac13,\tfrac56\}$ are also joined; (4) leaf diameters tend to $0$.
The quotients are $K'_B=\mathbb D/L_B$ and $J'_B=S^1/L_B$.  The depth of the leaf
$\langle\frac{a}{3\cdot2^n},\frac{b}{3\cdot2^n}\rangle$ is $n$.
\end{definition}

\begin{definition}[$z^2$-invariant lamination]
\label{def:matings-invariant}
$L_B$ is $z^2$-invariant if its leaves are simple curves and every leaf other than
$\langle\frac13,\frac23\rangle$ is an iterated preimage of $\langle\frac56,\frac16\rangle$ under $z\mapsto z^2$.
\end{definition}

\begin{definition}[Parameter lamination and mating]
\label{def:matings-parameter}
Using the parameter B\"ottcher map $B$, the parameter lamination is
$\overline L_B=\bigcup_{\ell\in L_B}\bigl((1/z)\circ B\bigr)^{-1}(\ell\cap\mathbb D)$.  Let $W_B$ be the
closed disk bounded by $\langle\frac13,\frac23\rangle$ and containing the Basilica limb.  Distinct
$c_1,c_2$ are equivalent if both lie in $W_B$, or lie on the same parameter leaf.  The mating is
$\M\cup_{S^1}L_B=(\M\cup_{S^1}\mathbb D)/L_B=V_1/L_B$, with quotient map $m\colon V_1\to V'_2$ and
$M'_2=m(\M)$.
\end{definition}

\begin{remark}
The paper does not define general laminations, quadratic lamination spaces, critical chords,
polygons, or gaps in the abstract sense; its laminations are the Basilica lamination and its leaves.
\end{remark}

\subsection*{Results}

\begin{theorem}[Dudko, Theorem~1.1]
\label{thm:matings-1-1}
The restriction $m|_\M\colon\M\to M'_2$ is canonical, independent of the choice of leaves.  $M'_2$ is
locally connected, and the quotients $V'_2(L_B)$ are homeomorphic, by homeomorphisms that preserve $M'_2$.
Status: proved in the paper (refined as its Proposition~3.8).
\end{theorem}

\begin{theorem}[Dudko, Theorem~1.2]
\label{thm:matings-1-2}
$\M\cup_{S^1}L_B$ is homeomorphic to $V_2$, preserving parapuzzle combinatorics and agreeing with the
straightening maps on every small Mandelbrot copy; the homeomorphism is unique over $M_2$.  Status: proved.
\end{theorem}

\begin{corollary}[Dudko, Corollary~1.3]
\label{cor:matings-1-3}
$M_2$ is locally connected and is a continuous image of $\M$.  Status: proved.
\end{corollary}

\begin{theorem}[Dudko, Theorem~1.4]
\label{thm:matings-1-4}
If $c\in\M\setminus W_B$ is not on the boundary of a hyperbolic component and is not four-times
renormalizable, then $f_c\cup_{S^1}L_B$ is well defined and topologically conjugate to $g_{m(c)}$.  If
$J_c$ is locally connected, then $f_c\cup_{S^1}L_B=f_c\cup_{S^1}f_B$.  Status: proved.
\end{theorem}

\begin{proposition}[Dudko, Propositions~3.2--3.4 and~3.8]
\label{prop:matings-3}
A small Mandelbrot copy $\M_H\not\subset W_B$ is bounded on the left (respectively right) by $L_B$
exactly when some $n$ satisfies $\frac56<2^n\phi_1<\frac16$ (respectively $\frac56<2^n\phi_2<\frac16$).  Every
such copy is bounded on at least one side, and satellite copies on both sides.  There is a generalized
decoration tiling of $\M_H$ in which each component of $V_1\setminus\widehat\M$ and each leaf of
$L_B\setminus W_B$ meets at most three pieces.  The boundaries of the components of $V_1\setminus\widehat\M$
contain no infinitely renormalizable parameter; these boundaries and $\widehat\M$ are locally connected;
the diameters of the complementary components tend to $0$.  Status: proved.
\end{proposition}

\subsection*{Relation to MLC}

The original mating construction of Wittner requires local connectivity of $\partial\M$.  The paper's
replacement by $\M\cup_{S^1}L_B$ does not.  Its principal consequence is Corollary~\ref{cor:matings-1-3}:
$M_2$ is locally connected.  The paper does not prove local connectivity of $\M$.

\subsection*{Dependencies}

The paper uses the Douady--Hubbard parameter B\"ottcher isomorphism $\C\setminus\M\cong\C\setminus\overline{\mathbb D}$
and the correspondence theorem for dynamic and parameter rays; Moore's quotient theorem; the extension of
homeomorphisms between locally connected planar continua; Yoccoz's puzzle and fiber results, the decoration theorem,
and Thurston rigidity; and earlier Basilica mating results of several authors.

\section{Core entropy}
\label{sec:digest-core-entropy}

\subsection*{Definitions}

\begin{definition}[Lamination core entropy]
\label{def:core-lamination}
For the invariant quadratic lamination $L_\vartheta$ let $N_\vartheta(n)$ count the generation-$n$
precritical leaves that lie in the backward orbit of the preferred major leaf and separate the
$\alpha$-gap from its negative.  Then $h(\vartheta)=\limsup_{n}\frac1n\log N_\vartheta(n)$.  If the
$\alpha$-gap is an infinite polygon, $h(\vartheta)=0$.
\end{definition}

\begin{definition}[Core entropy of a quadratic polynomial]
\label{def:core-polynomial}
For $c\in\M$ the critical ray pair is $\operatorname{RP}(\vartheta/2,(1+\vartheta)/2)$, where $\vartheta$ is an
external parameter angle.  A generation-$n$ precritical ray pair is a backward iterate taking $n$ iterates to the
critical ray pair (or to the characteristic periodic ray pair in the Fatou case).  If $\widetilde N_c(n)$ counts those
that separate the $\alpha$-fixed point from its negative, then $\widetilde h(c)=\limsup_n\frac1n\log\widetilde N_c(n)$.
\end{definition}

\begin{definition}[Entropy on $S^1$ and on $\M$]
\label{def:core-functions}
$h\colon S^1\to[0,\log2]$ is the lamination entropy, and $\widetilde h\colon\M\to[0,\log2]$ is obtained from the
projection $\M\to\M_{\rm abs}$ to the abstract Mandelbrot set, so that $\widetilde h=h\circ\pi$.  If the parameter ray at
angle $\vartheta$ lands or accumulates at $c$, then $h(\vartheta)=\widetilde h(c)$.
\end{definition}

\begin{definition}[Hubbard tree and topological entropy]
\label{def:core-hubbard}
For a postcritically finite polynomial, the Hubbard tree $H$ is the minimal tree in the filled Julia set connecting
the postcritical set.  If $M_{ij}\in\{0,1,2\}$ counts how often the edge $e_i$ covers $e_j$, then the classical core
entropy is $h_{\rm top}(p_c|H)=\log\lambda(M)$, where $\lambda(M)$ is the leading eigenvalue.
\end{definition}

\subsection*{Results}

\begin{theorem}[Dudko--Schleicher, Theorem~1.1; continuity]
\label{thm:core-1-1}
Both $h\colon S^1\to[0,\log2]$ and $\widetilde h\colon\M\to[0,\log2]$ are continuous.  Status: proved in the paper.
\end{theorem}

\begin{theorem}[Dudko--Schleicher, Theorem~1.2 and~4.1; maxima]
\label{thm:core-4-1}
As stated in the paper, and proved there: (1) every dyadic angle is an isolated local maximum of $h$;
(2) every local maximum is dyadic; (3) in every wake the unique global maximum occurs at the dyadic angle with
least denominator; (4) for each $n$, $N_\vartheta(n)$ is maximized in a wake at the dyadic of least generation.
These statements are the conjectures of Tiozzo in the form given by the paper.
\end{theorem}

\begin{corollary}[Dudko--Schleicher, Corollary~4.2]
\label{cor:core-4-2}
$\widetilde h$ has isolated local maxima exactly at the parameters $c(k/2^q)$.  Status: proved in the paper.
\end{corollary}

\begin{theorem}[Dudko--Schleicher, Theorem~4.3; minima]
\label{thm:core-4-3}
$h$ has an isolated local minimum at $\vartheta$ exactly when three rational parameter rays land at one
Misiurewicz--Thurston parameter with equal odd denominators; $\widetilde h$ has no isolated local minimum.
Status: proved in the paper; the printed inequality in the plain-text version was partly corrupted, so the
statement follows the displayed theorem formulation.
\end{theorem}

\begin{theorem}[Dudko--Schleicher, Theorem~2.4; correspondence]
\label{thm:core-2-4}
A rational dynamical ray pair is characteristic if and only if the corresponding parameter rays land together
and separate $c$ from $0$.  Status: proved.
\end{theorem}

\begin{theorem}[Dudko--Schleicher, Lemma~2.5; monotonicity]
\label{thm:core-2-5}
If $c\prec\widetilde c$, then $N_c(n)\le N_{\widetilde c}(n)$ for all $n$, and $h(c)\le h(\widetilde c)$.  Status: proved.
\end{theorem}

\begin{theorem}[Dudko--Schleicher, Theorems~5.7--5.9; continuity along veins]
\label{thm:core-5}
Entropy is continuous at hyperbolic-component angles; continuity along a vein implies continuity in a neighbourhood of
each point of the vein; hence $h$ is continuous everywhere.  Status: proved in the paper.
\end{theorem}

\begin{theorem}[Dudko--Schleicher, Theorem~A.3; biaccessibility]
\label{thm:core-a-3}
$\widetilde h(c)=B_{\rm top}(c)\log2$ and $h(\vartheta)=B_{\rm comb}(\vartheta)\log2$, where $B_{\rm top}$ and $B_{\rm comb}$ are
the topological and combinatorial biaccessibility dimensions.  Corollary~A.5: $h_{\rm top}(T_c)=\widetilde h(c)$ whenever the
tree $T_c$ is compact.  Status: proved in the paper.
\end{theorem}

\subsection*{Relation to MLC}

The paper does not prove local connectivity of $\M$.  It treats parameters whose rays need not land, and it works in the
abstract Mandelbrot set to avoid assuming MLC.  The entropy function extends continuously to all parameters.  Positive
entropy is related to biaccessibility dimension (Theorem~\ref{thm:core-a-3}); the paper does not assert that positive
entropy implies MLC.  The paper records known consequences of Yoccoz puzzles along dyadic veins and treats general MLC as open.

\subsection*{Dependencies}

Thurston's invariant laminations and correspondence theorem; Milnor's parameter rays; the Douady--Hubbard branch theorem and
pinched-disk model; Yoccoz's theorem and puzzles; Tiozzo's continuity results and conjectures; Jung's appendix; Bruin--Schleicher's
dimension results; and Furstenberg's relation between entropy and dimension.

\section{Holomorphic motions}
\label{sec:digest-slodkowski}

Słodkowski's paper~\cite{Slodkowski} does not mention the Mandelbrot set.  It is included because the
holomorphic-motion arguments of the phase--parameter approach rely on it (Chapter~\ref{ch:phase-parameter}).

\subsection*{Definitions}

\begin{definition}[Holomorphic motion]
\label{def:holomorphic-motion}
Let $E\subset\C$.  A holomorphic motion of $E$ parametrized by the unit disk $\D$ is a map $f\colon\D\times E\to\C$ such that:
(a) for each $w\in E$, $z\mapsto f(z,w)$ is holomorphic on $\D$; (b) for each $z\in\D$, $w\mapsto f_z(w)=f(z,w)$ is injective; (c) $f_0$ is the identity on $E$.
\end{definition}

\begin{definition}[Analytic disk, trace, polynomial hull]
\label{def:holomorphic-trace}
An analytic disk is the graph of a holomorphic map $\D\to\C$.  The trace of a motion is
$\{(z,f_z(w)):z\in\D,\ w\in E\}\subset\D\times\C$.  The polynomial hull $h(X)$ of a compact set $X$ is the usual polynomially convex hull.
\end{definition}

\subsection*{Results}

\begin{lemma}[Lambda lemma; cited from Mañé--Sad--Sullivan]
\label{lem:lambda-lemma}
A holomorphic motion is jointly continuous and extends continuously to $\overline{\D}\times\overline E$.  Its injections
$F_z$ are quasiconformal.  Status: cited.
\end{lemma}

\begin{theorem}[Słodkowski, Theorem~1.2; cited from the author's earlier work]
\label{thm:slodkowski-1-2}
Let $X\subset\partial\D\times\C$ be compact with connected fibres $X(\zeta)$ for all $\zeta\in\partial\D$, and suppose
$h(X)\cap(\D\times\C)\neq\varnothing$.  Then $h(X)\setminus X$ is a union of analytic disks, and the topological boundary of
$h(X)\setminus X$ in $\D\times\C$ is a union of mutually disjoint analytic disks.  Status: cited.
\end{theorem}

\begin{theorem}[Słodkowski, Theorem~1.3; extension]
\label{thm:slodkowski-1-3}
Every holomorphic motion $f\colon\D\times E\to\C$ of an arbitrary subset $E\subset\C$ extends to a holomorphic motion
$F\colon\D\times\C\to\C$ of all of $\C$.  Status: proved in the paper.  The earlier results of Sullivan--Thurston and of
Bers--Royden gave extensions over smaller disks, with radius $r$ unspecified and $r=3$ respectively; the paper obtains
$r=1$.
\end{theorem}

\begin{theorem}[Słodkowski, Theorem~1.4; holomorphic axiom of choice]
\label{thm:slodkowski-1-4}
Let $f(z,w)=f_z(w)$ be a holomorphic motion of $E\subset\C$ parametrized by $z\in\D$.  For every point $a\notin E$ there is a
holomorphic map $g\colon\D\to\C$ with $g(0)=a$ and $g(z)\notin f_z(E)$ for every $z\in\D$.  Status: proved in the paper for
finite sets (its Section~2), and sketched in general in its Section~3.  The paper notes that the axiom of choice gives the
extension theorem by a known argument of Sullivan and Thurston.
\end{theorem}

\subsection*{Relation to MLC}

The text contains no statement about the Mandelbrot set.  Its role for MLC is technical: it supplies the extension of holomorphic motions
used when the phase--parameter correspondence is built from the trace of a motion, which the review treats in
Chapter~\ref{ch:phase-parameter}.
